\cotefile{157-1}
\scanpage{82}
\mspage{\unsure{88|28}}

On pose

\[
X'=S'\times\overline S\qquad X''=S''\times\overline S
\]

\margin{$\overline S$ copie conforme de $S$, cf ci-dessus}

et il faut, pour recoller $X'$ et $X''$, se donner un \unsure{isom|iso} entre $X'\mid\unsure{\Omega|R}\simeq\unsure{\mathbb{T}|T|R|\mathbb{I}}\times\overline S=\unsure{\overline S'|\overline S}\amalg\overline S'$ et $X''\mid\Omega\simeq\mathbb{T}\times\overline S=\unsure{\overline S''|\overline S}\amalg\overline S''$ qui \deleted{\unsure{respecte la}} soit un $T$-foncteur, i.e. qui provienne de deux iso.

\diagram{\begin{tikzcd}[ampersand replacement=\&,row sep=small]
\overline S\arrow[r,"u_a"] \&\overline S\\
\overline S\arrow[r,"u_b",no head]\&\overline S .
\end{tikzcd}}

On prendra \deleted{\unsure{donc|$u_a=$}}

\[
\begin{aligned}
u_a&=\mathrm{id}_{\overline S}\\
u_b&=\text{symétrie p.r. à \unsure{l'axe} $\overline0,\overline{0'}$ i.e.,}\\
&\qquad u_b\text{ laisse fixes }\overline0,\overline{0'}\text{ et intervertit}\\
&\qquad\overline a\text{ et }\overline b .
\end{aligned}
\]

(C'\deleted{\illegible{1}}\inserted{est} plus clair si on \unsure{désidentifie|désidentifiait} la base et la fibre, en prenant fibre $\overline S$ :

\diagram{\begin{tikzcd}[ampersand replacement=\&,row sep=tiny,column sep=small]
\&\overline a\&\\
\overline0\arrow[ur]\arrow[dr]\&\&\overline{0'}\arrow[ul]\arrow[dl]\\
\&\overline b\&
\end{tikzcd}}

)

\deleted{\unsure{Le|Il}} Pour voir que $X$ est une \deleted{\unsure{$\mathscr F$|$f$}} $\mathbb D$-fibration sur $S$ (quel que soit $\mathbb D$) \deleted{en fait il} \inserted{\unsure{prouvons}} \deleted{il suffit de voir} qu'il \unsure{est} \uline{propre} sur $S$ --- \deleted{ce qui \inserted{est}} \unsure{revient|revient à} : voir que $X'\to S'$ et $X''\to S''$ sont propres, ce qui est trivial. \deleted{Il suffit} Pour voir que c'est une \inserted{$\mathbb D$-}fibration, on est également ramené au cas de $X'\to \unsure{S'|O'}$ et $X''\to S''$, où c'est aussi trivial. On a utilisé le

% slanted
\margin{\uline{NB} $f:X\to S$ est également \uline{lisse}, car c'est un \unsure{fibré trivial} (un produit) au-dessus des deux \deleted{\unsure{$+$}} ouverts $S\smallsetminus\{a\}$, $S\smallsetminus\{b\}$, qui recouvrent $S$. Mieux : $f$ est à la fois $\mathit{Cat}$-fibrant et $\mathit{Cat}$-cofibrant (Caractérisation \unsure{infinitésimale} par les propriétés des $X\times_S\Delta^{\unsure{2|n}}\to\Delta^2$, pour $\Delta^2\to S$).}

\uline{Lemme} Soit $f:X\to S$, $(S_i)_{i\in I}$ une famille de sous-\deleted{catégories} fermés recouvrant $S$. $\forall i$, soit $X_i=X_i\times_S S_i\xrightarrow{f_i}S_i$. Pour que $f$ soit propre, il faut et que les $f_i$ le soient. \deleted{Pour que $f$ soit une} $\mathbb D$-fibration propre, il f. et s. que les $f_i$ le soient.

% archive sheet number: 81
